% Reconstructed from the compiled lecture-note PDF.
\chapter{Hamiltonian mechanics and optimal control}
This appendix explains the vocabulary of cotangent bundles, Hamiltonians, costates, and the Pontryagin maximum principle. The exposition begins in Euclidean space and then indicates what changes on a matrix group.


\section{Tangent and cotangent vectors}

\begin{displaytext}
For an open set M ⊂Rn, a tangent vector at q ∈M is an ordinary vector v ∈Rn. A cotangent \
vector is a linear functional on tangent vectors. In coordinates it is represented by a row vector or \
by a one-form \
p = p1 dq1 + · · · + pn dqn.
\end{displaytext}

\begin{displaytext}
The pairing is \
n \
⟨p, v⟩∗= X pjvj. \
j=1 \
The collection of all pairs (q, p) is the cotangent bundle T ∗M. \
A smooth manifold looks locally like an open set in Rn, so tangent and cotangent vectors are defined \
by changing coordinates consistently. The cotangent bundle always has twice the dimension of M.
\end{displaytext}

\section{The canonical symplectic form}

On T ∗Rn with coordinates (q1, . . . , qn, p1, . . . , pn), define the canonical one-form


\begin{displaytext}
n \
θ = X pj dqj \
j=1
\end{displaytext}

\begin{displaytext}
and symplectic form \
n \
ω = −dθ = X dqj ∧dpj. \
j=1
\end{displaytext}

\begin{displaytext}
A two-form is symplectic if it is closed and nondegenerate. Nondegeneracy means that for every \
differential dH there is a unique vector field H satisfying \
ω(H, ·) = dH. \
This H is the Hamiltonian vector field of H.
\end{displaytext}

\begin{displaytext}
In coordinates the equation becomes Hamilton’s equations: \
∂H q′j = , p′j = −∂H . \
∂pj ∂qj \
Example C.1 (Ordinary mechanics). For \
||p||2 \
H(q, p) = + V (q), \
2m \
Hamilton’s equations give \
p \
q′ = m, p′ = −∇V (q), \
which combine to Newton’s equation mq′′ = −∇V (q).
\end{displaytext}

\section{Poisson brackets}

For smooth functions F, G on phase space, define


\begin{displaytext}
n ∂F ∂G ∂G ! \{F, G\} = X −∂F . \
∂qj ∂pj ∂pj ∂qj j=1
\end{displaytext}

\begin{displaytext}
Along the Hamiltonian flow of H, \
d \
dtF(q(t), p(t)) = \{F, H\}. \
Therefore F is conserved if \{F, H\} = 0. The bracket is bilinear, antisymmetric, satisfies a product \
rule, and obeys the Jacobi identity.
\end{displaytext}

The Fuller system in the source uses a nonstandard but explicit symplectic form, and therefore a nonstandard Poisson bracket. The principle is unchanged.


\section{From a control problem to a Hamiltonian}

\begin{displaytext}
Consider \
q′ = f(q, u), u(t) ∈U,
\end{displaytext}

\begin{displaytext}
with running cost \
Z tf \
L(q(t), u(t)) dt. \
0 \
Introduce a costate p(t) ∈T q(t)M∗ and a constant multiplier λ0 ≤0. Define the control-dependent \
Hamiltonian \
H(q, p, u) = ⟨p, f(q, u)⟩∗+ λ0L(q, u). \
The first term rewards motion aligned with the costate, while the second incorporates the objective.
\end{displaytext}

\begin{displaytext}
The sign convention varies across textbooks. The source uses maximization with λ0 ≤0. In the \
normal case one rescales to λ0 = −1.
\end{displaytext}

\section{The Pontryagin maximum principle}

Pontryagin maximum principle, free terminal time


\begin{displaytext}
Let (q∗, u∗) be an optimal trajectory for a sufficiently regular finite-dimensional control \
problem. Then there exist an absolutely continuous costate p(t) and a constant λ0 ≤0, not \
both zero, such that almost everywhere:
\end{displaytext}

\begin{displaytext}
∂H \
q′ = = f(q, u∗), \
∂p \
p′ = −∂H , \
∂q \
H(q, p, u∗) = max H(q, p, u). \
u∈U
\end{displaytext}

For an autonomous problem with free terminal time, the maximized Hamiltonian is constant and equals zero along the extremal. Endpoint constraints produce transversality conditions on p.


A trajectory satisfying these conditions is a Pontryagin extremal. The principle gives necessary, not sufficient, conditions for optimality.


\subsection{Normal and abnormal extremals}

\begin{displaytext}
If λ0 < 0, rescale all multipliers and set λ0 = −1. This is a normal extremal. If λ0 = 0, the \
cost disappears from the Hamiltonian; the extremal is abnormal. Abnormality can reflect genuine \
geometry or an artifact of how constraints were encoded.
\end{displaytext}

\subsection{Why the Hamiltonian is maximized}

A short needle variation changes the control from u∗to v on a tiny interval. The first-order change of endpoint and cost is measured by the Hamiltonian difference


H(q, p, v) −H(q, p, u∗).


Optimality forces this difference to be nonpositive for every v, hence maximization by u∗.


\section{Bang-bang controls}

Suppose U is a compact convex polytope and H depends affinely on u. A nonconstant affine function on a polytope reaches its maximum on a face, usually at a vertex. Thus optimal controls often take extreme values.


In the Reinhardt system the control-dependent term is a ratio of affine functions,


−⟨ΛR, Zu⟩ ⟨X, Zu⟩,


but along every segment of the control triangle it is monotone. Hence the maximizing set is still a face. On the open regions where a unique vertex wins, the control is bang-bang.


\begin{displaytext}
A switching function compares two candidate vertices. If controls i and j give Hamiltonians Hi, Hj, \
define \
χij = Hi −Hj.
\end{displaytext}

\begin{displaytext}
The wall χij = 0 is where a switch may occur. A transverse zero gives an isolated switch. An \
infinite sequence of switches in finite time is chattering.
\end{displaytext}

\section{Singular controls}

A singular interval is one on which the maximum condition fails to determine a unique control. One differentiates switching functions until the control appears, then solves the resulting equations. This procedure may identify a special interior control.


For the Reinhardt problem, full indeterminacy forces


\begin{displaytext}
3 \
ΛR = 0, X = J, Λ1 = 2λ0J.
\end{displaytext}

The corresponding control is the center of the triangle,


u = (1/3, 1/3, 1/3),


and the boundary curve is circular. Thus the singular locus has a concrete geometric meaning.


\section{Transversality}

\begin{displaytext}
Suppose the final state is constrained to lie on a submanifold N ⊂M. A permissible endpoint \
variation v lies in Tq(tf)N. The boundary term in the first variation is ⟨p(tf), v⟩∗. Optimality \
requires \
⟨p(tf), v⟩∗= 0 for all v ∈Tq(tf)N.
\end{displaytext}

Thus the costate is normal to the endpoint manifold.


In the Reinhardt problem, the endpoint is periodic modulo the rotation R. Consequently the state and costate satisfy conjugacy relations such as


\begin{displaytext}
X(tf) = R−1X(0)R, ΛR(tf) = R−1ΛR(0)R.
\end{displaytext}

These conditions allow the trajectory to be extended periodically after applying the discrete symmetry.


\section{Left-invariant systems on a matrix group}

For g ∈SL2(R), write the state equation


g′ = gX.


\begin{displaytext}
The vector field is left-invariant because replacing g by hg changes the velocity to hgX. A cotangent \
vector at g can be transported to the identity and represented by a matrix Λ1 ∈sl2(R) using the \
trace pairing.
\end{displaytext}

The Hamiltonian contribution from the group equation is then


⟨Λ1, X⟩.


Hamilton’s equation for the left-trivialized costate becomes


\begin{displaytext}
Λ′1 = [Λ1, X].
\end{displaytext}

\begin{displaytext}
Its solution is \
Λ1(t) = g(t)−1Λ1(0)g(t),
\end{displaytext}

so det Λ1 is conserved.


\section{The Reinhardt Hamiltonian}

\begin{displaytext}
The state is (g, X) and the reduced costates are (Λ1, ΛR). With the trace pairing, the Hamiltonian \
is \
Zu⟩ H(Λ1, ΛR, X; u) = Λ1 −3 X −⟨ΛR, 2λ0J, ⟨X, Zu⟩. \
Set \
Zu \
P = ⟨X, Zu⟩. \
The state–costate equations are
\end{displaytext}

\begin{displaytext}
g′ = gX, \
X′ = [P, X], \
Λ′1 = [Λ1, X], \
Λ′R = [P, ΛR] −⟨ΛR, P⟩[P, X] + [−Λ1 + 32λ0J, X].
\end{displaytext}

Every term is a 2 × 2 commutator or trace pairing. The system looks formidable, but it has three simplifying features:


\begin{displaytext}
• X remains on a two-dimensional adjoint orbit; \
• ΛR is constrained by ⟨ΛR, X⟩= 0; \
• the control triangle has only three vertices and three walls.
\end{displaytext}

\section{Noether’s principle and angular momentum}

In mechanics, a continuous symmetry of the Hamiltonian produces a conserved quantity. This is Noether’s theorem. A version remains valid for optimal control.


\begin{displaytext}
For circular control, the control set is invariant under all rotations. Simultaneous rotation of \
X, Λ1, ΛR, and the control leaves the Hamiltonian unchanged. The corresponding momentum is
\end{displaytext}

\begin{displaytext}
A = ⟨J, Λ1 + ΛR⟩.
\end{displaytext}

It is constant along circular-control extremals.


The triangular control set has only discrete rotational symmetry, so this continuous conservation law is lost. The circular problem is therefore a symmetry-enhanced model used to discover the Fuller normal form.


\section{Second variation}

The maximum principle detects only first-order obstructions. A Pontryagin extremal can fail to minimize because of a negative second variation.


\begin{displaytext}
For a one-parameter family qs with q0 = q, write \
s2 \
J (qs) = J (q0) + s δJ + δ2J + o(s2). \
2 \
At an extremal, δJ = 0. If one can choose a variation with δ2J < 0, then q is not a local minimizer.
\end{displaytext}

The source applies this to the circular singular arc. A compactly supported matrix perturbation has second-order area term proportional to


Z − (ψ′1ψ2 −ψ1ψ′2) dt.


Choosing the signs of ψ′1, ψ′2 appropriately makes this negative while preserving positive curvature for sufficiently small perturbations. Thus the circle is a Pontryagin extremal but not minimizing.


\section{Hamilton–Jacobi viewpoint}

For intuition, define the value function V (q) to be the least cost from q to the target. Formally, V satisfies a Hamilton–Jacobi–Bellman equation


\begin{displaytext}
0 = min L(q, u) + dVq(f(q, u)) . \
u∈U
\end{displaytext}

The costate along an optimal trajectory is related to dV . A globally smooth value function rarely exists in switching problems, but this viewpoint explains why a sharp global subsolution could prove the full Reinhardt conjecture without classifying every extremal.


\section{Exercises}

\begin{statementbox}{Exercise}
Exercise C.2. For H(q, p) = pf(q) −L(q) in one dimension, write Hamilton’s equations and verify directly that H is constant when it has no explicit time dependence.
\end{statementbox}

\begin{statementbox}{Exercise}
Exercise C.3. Let U = [−1, 1] and H = pu. Determine the maximizing control. What happens when p = 0?
\end{statementbox}

\begin{statementbox}{Exercise}
Exercise C.4. Suppose Hi −Hj has a simple zero at t0. Explain why the maximizing control switches only once near t0. Exercise C.5. Verify that Λ1(t) = g(t)−1Λ1(0)g(t) solves Λ′1 = [Λ1, X] when g′ = gX.
\end{statementbox}

\begin{insightbox}{Chapter summary}
\end{insightbox}

Pontryagin’s principle lifts the state trajectory to a Hamiltonian trajectory with costates. The control maximizes a pointwise Hamiltonian, producing vertex controls, switching walls, and singular loci. In the Reinhardt problem, left invariance turns the abstract equations into explicit commutator identities on 2 × 2 matrices.

